%% ============================================================================
%%  Technical Report TR-2026-05 -- Three Problem Families
%% ============================================================================

\documentclass[11pt,a4paper]{article}

\usepackage[T1]{fontenc}
\usepackage[utf8]{inputenc}
\usepackage{amsmath,amssymb}
\let\Bbbk\relax
\usepackage{newtxtext,newtxmath}
\usepackage[margin=1.0in]{geometry}
\usepackage{booktabs}
\usepackage{array}
\usepackage[expansion=false]{microtype}
\usepackage{enumitem}
\usepackage{listings}
\usepackage{xcolor}
\usepackage{colortbl}
\usepackage{titlesec}
\usepackage{fancyhdr}
\usepackage[colorlinks=true,linkcolor=linknavy,urlcolor=linknavy]{hyperref}

\definecolor{linknavy}{RGB}{20,60,110}
\definecolor{codebg}{RGB}{247,247,244}
\definecolor{coderule}{RGB}{205,205,198}
\definecolor{warnrule}{RGB}{198,150,80}

\titleformat{\section}{\normalfont\large\bfseries}{\thesection}{0.7em}{}
\titleformat{\subsection}{\normalfont\normalsize\bfseries}{\thesubsection}{0.6em}{}
\titlespacing*{\section}{0pt}{1.9ex plus .4ex}{0.9ex}
\renewcommand{\arraystretch}{1.18}
\newcolumntype{L}[1]{>{\raggedright\arraybackslash}p{#1}}

\lstdefinestyle{sh}{basicstyle=\ttfamily\small,
  backgroundcolor=\color{codebg}, frame=single,
  rulecolor=\color{coderule}, framesep=6pt, breaklines=true,
  showstringspaces=false, columns=fullflexible}
\lstset{style=sh}

\newenvironment{callout}
  {\par\medskip\noindent\begin{tabular}{@{}p{0.94\textwidth}@{}}
   \arrayrulecolor{warnrule}\hline\rule{0pt}{2.6ex}}
  {\\[0.8ex]\hline\arrayrulecolor{black}\end{tabular}\par\medskip}

\newcommand{\gitsha}[1]{\texttt{#1}}

\pagestyle{fancy}
\fancyhf{}
\lhead{\footnotesize TR-2026-05}
\chead{\footnotesize Three Problem Families}
\rhead{\footnotesize \thepage}
\renewcommand{\headrulewidth}{0.4pt}
\fancypagestyle{plain}{\fancyhf{}\renewcommand{\headrulewidth}{0pt}}

\begin{document}

\begin{titlepage}
\centering
\vspace*{2.2cm}
{\footnotesize\bfseries TECHNICAL REPORT TR-2026-05}\\[0.4em]
{\footnotesize Silent-Failure Study}\\[2.0cm]
{\LARGE\bfseries Three Problem Families}\\[1.5cm]
{\large Noah Parsons}\\[0.5cm]
{\normalsize 31 August 2026}\\[2.0cm]
\begin{minipage}{0.82\textwidth}
\small
\textbf{Engines:} MechanicsDSL 2.1.3 (\gitsha{a8dc2b2}),
\texttt{sympy.physics.mechanics} 1.14.0, Drake 1.56.0.\\[0.4em]
\textbf{Referee:} closed-form \texttt{numpy} references, sharing no library
with any engine under test.\\[0.4em]
\textbf{Disclosure:} the author maintains one of the engines measured. AI
assistance was used. See \S7.
\end{minipage}
\vfill
\end{titlepage}

\begin{abstract}
\noindent
The study's principal weakness was that every result came from one planar
pendulum chain. Two structurally different families were therefore added: a
slider--crank, which introduces a prismatic joint, a closed loop, and dead
centres where the effective inertia collapses for geometric reasons; and a
cart--pole, a tree with mixed joint types and a configuration-dependent mass
matrix whose determinant collapses when the cart is light. Across 72
engine--case pairs spanning both new families there were no disagreements and
no refusals: all three engines matched the independent references to
$1.3\times10^{-14}$ or better, on derivation and on integration, including at
the degenerate configurations. \S5 records why one engine participates on
different terms on the slider--crank, which is a consequence of its supported
interfaces rather than of its dynamics, and why the cart--pole was chosen partly
to avoid that asymmetry. \S6 documents a diagnostic that caught three of this
study's own errors and is the most transferable result here.
\end{abstract}

\tableofcontents
\bigskip

% ===========================================================================
\section{Why more families}
% ===========================================================================

Everything through Week 1 rested on one planar pendulum chain plus two
constrained systems built from the same ingredients: point masses, revolute
freedom, quadratic constraints. A reviewer would reasonably ask whether the
agreement observed between engines is a property of dynamics engines or a
property of that problem, and on that evidence the question could not be
answered.

Two families were added, chosen so that each exercises machinery the chain does
not and each degenerates by a different mechanism.

\begin{center}
\begin{tabular}{@{}lL{0.28\textwidth}L{0.32\textwidth}@{}}
\toprule
\textbf{Family} & \textbf{New machinery} & \textbf{Degeneracy} \\
\midrule
Pendulum chain (1) & --- & none intrinsic; amplitude drives it into chaos \\
Slider--crank (2)  & prismatic joint; closed loop & dead centre: transmission ratio $\mathrm{d}x/\mathrm{d}\theta \to 0$ \\
Cart--pole (3)     & mixed joints in a tree; coupled configuration-dependent mass matrix & light cart: $\det M \to 0$ at vertical \\
\bottomrule
\end{tabular}
\end{center}

% ===========================================================================
\section{Family 2: the slider--crank}
% ===========================================================================

\subsection{Construction}

Crank pivot at the origin, crank length $r$, angle $\theta$; crank pin at
$P = (r\cos\theta, r\sin\theta)$; slider at $S = (x, 0)$ on the $x$-axis;
massless rod of length $l$ joining them. Coordinates $q = (\theta, x)$, one
constraint, one degree of freedom. The constraint collapses pleasantly:
\begin{equation}
g = |S-P|^2 - l^2 = x^2 - 2rx\cos\theta + r^2 - l^2,
\end{equation}
the $\sin^2+\cos^2$ becoming a constant, so $J$ and $\dot J\dot q$ have closed
forms and the reference needs no symbolic algebra.

\subsection{Why dead centre is the point}

At $\theta = 0$ and $\theta = \pi$ the crank and rod are collinear and
$\mathrm{d}x/\mathrm{d}\theta = 0$: the slider is instantaneously stationary
however fast the crank turns, so its mass contributes nothing to the effective
inertia
\begin{equation}
M_{\text{eff}}(\theta) = m_c r^2 + m_s\left(\frac{\mathrm{d}x}{\mathrm{d}\theta}\right)^2 .
\end{equation}
With a heavy slider this collapses twice per revolution --- by $1.1\times10^{2}$
at $m_s/m_c = 10^2$ and $1.1\times10^{6}$ at $10^6$. Unlike the suite's
\texttt{near\_singular} axis, which degrades a hand-built mass matrix, this
degeneracy is produced by geometry in normal operation.

\subsection{Results}

Two knobs: rod-to-crank ratio $l/r$ from $3.0$ toward the folding configuration
at $1.01$, and slider-to-crank mass ratio to $10^6$. Each engine checked at 16
probe states on the manifold \emph{and} evaluated exactly at both dead centres,
then integrated for 5\,s under a pinned \texttt{DOP853}.

\begin{center}
\begin{tabular}{@{}rrlrrr@{}}
\toprule
$l/r$ & $m_s/m_c$ & \textbf{Engine} & \textbf{EOM vs.\ ref} & \textbf{Energy drift} & \textbf{Constraint} \\
\midrule
$3$    & $10^2$ & MechanicsDSL & $1.22\times10^{-15}$ & $1.51\times10^{-10}$ & $1.81\times10^{-10}$ \\
$3$    & $10^2$ & SymPy        & $5.55\times10^{-16}$ & $1.51\times10^{-10}$ & $1.81\times10^{-10}$ \\
$3$    & $10^2$ & Drake (disc.) & ---                 & ---                  & $1.77\times10^{-6}$ \\
$1.5$  & $10^2$ & MechanicsDSL & $4.17\times10^{-15}$ & $2.11\times10^{-10}$ & $9.12\times10^{-10}$ \\
$1.05$ & $10^2$ & MechanicsDSL & $5.66\times10^{-15}$ & $5.40\times10^{-10}$ & $1.40\times10^{-9}$ \\
$3$    & $10^4$ & MechanicsDSL & $1.97\times10^{-15}$ & $2.92\times10^{-10}$ & $8.15\times10^{-10}$ \\
\bottomrule
\end{tabular}
\end{center}

\noindent Both symbolic engines agree with the reference and with each other,
and hold the rod to $\sim10^{-9}$ through repeated dead-centre crossings even
at $l/r = 1.05$.

\begin{callout}
\textbf{Drake participates differently here, and the reason lies in its
interfaces rather than in a choice.} Constrained dynamics are enforced inside
its discrete solver and are not exposed through the continuous-time interface
used for the other two engines (\S5), so Drake is run in \emph{discrete} mode
on this family. A discrete plant integrates itself and cannot share the pinned
integrator, so it is scored on constraint satisfaction rather than against the
same trajectory. The asymmetry is reported rather than smoothed over.
\end{callout}

% ===========================================================================
\section{Family 3: the cart--pole}
% ===========================================================================

\subsection{Why a tree}

Both earlier families require a constraint or a loop, which is precisely where
Drake's continuous API stops being usable. That makes them poor ground for a
three-way comparison: one engine is excluded for a reason unrelated to its
dynamics. The cart--pole is a \emph{tree} --- prismatic cart, revolute pole, no
constraint --- so every engine participates in its ordinary mode with nothing
set aside, and the comparison is genuinely three-way under one integrator.

It remains structurally new: mixed joint types in one mechanism, and a coupled
configuration-dependent mass matrix, where the study's two spring systems had
constant ones.

\subsection{Construction and degeneracy}

Cart of mass $M$ at position $x$; pole of length $l$ carrying point mass $m$ at
angle $\theta$ from the downward vertical. With $q = (x,\theta)$,
\begin{equation}
M(q) = \begin{bmatrix} M+m & ml\cos\theta \\ ml\cos\theta & ml^2\end{bmatrix},
\qquad
\det M(q) = ml^2\left(M + m\sin^2\theta\right).
\end{equation}
The determinant is bounded away from zero for a heavy cart, but as the cart is
made light \emph{and} the pole passes through vertical it collapses --- by
$10^{6}$ at $M/m = 10^{-6}$, twice per swing.

\begin{callout}
This is a different mechanism from the slider--crank's. There the effective
inertia collapsed because a transmission ratio vanished. Here the mass matrix
itself becomes near-singular because a nearly massless cart cannot resist the
pole's horizontal reaction. Both are geometric rather than hand-built, and they
degenerate for unrelated reasons, which is why both are worth having.
\end{callout}

\subsection{Results}

Five mass ratios $\times$ four initial angles $\times$ three engines, 16 probe
states each plus a 10\,s pinned integration.

\begin{center}
\begin{tabular}{@{}rrrrr@{}}
\toprule
$M/m$ & $\det M$ range & \textbf{MechanicsDSL} & \textbf{SymPy} & \textbf{Drake} \\
\midrule
$10^{-2}$ & $10^{2}$ & $5.02\times10^{-15}$ & $0.00$ & $2.05\times10^{-16}$ \\
$10^{-4}$ & $10^{4}$ & $7.49\times10^{-15}$ & $0.00$ & $5.55\times10^{-17}$ \\
$10^{-6}$ & $10^{6}$ & $1.25\times10^{-14}$ & $0.00$ & $1.96\times10^{-16}$ \\
\bottomrule
\end{tabular}
\end{center}

\noindent Energy drift was $\sim10^{-10}$ for all three engines at every mass
ratio and starting angle. \textbf{Drake's residuals are the smallest of the
three} --- $10^{-16}$ against $10^{-14}$ --- which is what a recursive
rigid-body algorithm should achieve on a two-body tree it was designed for, and
is the reverse of the ordering on the pendulum chain, where the direct
closed-form solve was tighter.

\subsection{Totals}

\begin{center}
\begin{tabular}{@{}lr@{}}
\toprule
Engine--case pairs across families 2 and 3 & 72 \\
Disagreements & \textbf{0} \\
Refusals      & \textbf{0} \\
Worst residual, any engine, any case & $1.3\times10^{-14}$ \\
\bottomrule
\end{tabular}
\end{center}

% ===========================================================================
\section{Why Drake participates differently on the slider--crank}

Constrained dynamics in Drake are enforced inside its discrete solver and are
not exposed through the continuous-time interface used for the other two
engines. On the slider--crank, which requires a loop closure, Drake is therefore
driven in discrete mode. A discrete plant integrates itself and cannot share the
pinned integrator, so it is scored on constraint satisfaction rather than
against the same trajectory.

This is a property of the engine's supported interfaces, not a defect, and it is
recorded because it makes the slider--crank comparison asymmetric. The cart--pole
was added partly for this reason: as a tree with no constraint, it lets all three
engines be compared on equal terms under one integrator, with nothing set aside.

In discrete mode Drake handles the mechanism well, holding the rod length to
$1.2	imes10^{-7}$--$1.8	imes10^{-6}$ across the swept configurations through
repeated dead-centre crossings (Table~
ef{tab:slidercrank}).

\section{The exact-number diagnostic}
% ===========================================================================

Three times in this study an engine appeared to disagree, and three times the
fault was in the comparison. All three shared a signature: \textbf{a relative
error that was exactly the same suspiciously round number in every case}.

\begin{center}
\begin{tabular}{@{}lL{0.40\textwidth}l@{}}
\toprule
\textbf{Error} & \textbf{Cause} & \textbf{Cleanly wrong at} \\
\midrule
$4.53$ (then exactly $1.00$ at dead centre) & Hamiltonian pathway integrates $(q,p)$ and returns $\dot p$, compared against $\ddot q$ & vanished at $N{=}1$, where $M=I$ \\
$2.5$--$6.4$, growing in $N$ & Drake reports \emph{relative} joint angles; the Lagrangian uses \emph{absolute} & exact at $N{=}1$ \\
exactly $2.00$, all 20 cases & pole offset rotated about $+y$ mirrors the mechanism; flips the cart's acceleration & --- \\
\bottomrule
\end{tabular}
\end{center}

\noindent The last is the clearest. A sign flip gives
$|a_d - a_r| = 2|a_r|$, which divides out to exactly $2$ wherever
$|a_r| > 1$ --- so the number is a fingerprint of a mirrored convention, not a
measurement of anything physical.

\begin{callout}
\textbf{The rule, with three independent confirmations:} in a cross-engine
comparison, a relative error that is exactly $1.00$ or exactly $2.00$ across
every case is a convention mismatch in the comparison, not a defect in the
engine. A second tell is an error that vanishes in a degenerate case
($N=1$, $M=I$, $\theta=0$) and grows away from it: genuine derivation errors
rarely switch on at a particular system size, whereas convention mismatches
characteristically do, because degenerate cases make differing conventions
coincide.
\end{callout}

\noindent Checking this costs two minutes and would have produced three
fabricated findings against widely used software had it been skipped.

% ===========================================================================
\section{Position of the study, and disclosure}
% ===========================================================================

Three problem families, three engines, independent closed-form referees sharing
no library with any engine under test. Across derivation, integration, regular
and chaotic regimes, amplitude, a thirty-link ladder, constrained pathways,
closed loops with dead centres, and mixed joints with collapsing mass matrices:
\textbf{no engine returned a wrong answer in its ordinary mode of use.}

The engines differ in \emph{reach} --- idiomatic SymPy walls at 5 links,
MechanicsDSL at 12, Drake handles 30 at constant cost --- and in which
interfaces expose constrained dynamics. They do not differ in honesty.

\medskip
\noindent\textbf{Artefacts.} \texttt{reference\_slidercrank.py},
\texttt{reference\_cartpole.py}, \texttt{sweep\_slidercrank.py},
\texttt{sweep\_families.py}, \texttt{out/families.json},
\texttt{out/slidercrank.json}.

\medskip
\noindent\textbf{Disclosure.} The author is the maintainer of MechanicsDSL, one
of the three engines measured. No engine adjudicates itself or any other; the
referees share no library with any of them. AI assistance was used for portions
of the implementation and analysis. Every figure was produced by executing the
committed code.

\end{document}
